\section{Introduction}
\label{sec:intro}

A dashcam records a rear-end collision: the camera car closes on a slow black truck and strikes it.
Asked to explain the crash, a strong open Video-LLM returns three paragraphs on overcast skies, road surface, and ``traffic conditions''.
The truck is never mentioned.
A domain-adapted version of the same model does name a lead vehicle and a rear-end impact---but calls the truck a red car, and places the impact in a one-second window it uses for more than half of all videos in the test set.

These two outputs fail differently, and the difference matters.
The first \emph{omits}: it says nothing false, but nothing useful.
The second \emph{hallucinates}: it commits to specific, checkable claims that are wrong.
For an insurer, an investigator, or an autonomous-driving safety team, the second failure is worse, because a confident and specific narrative is precisely the kind that gets acted upon~\cite{maynez2020faithfulness,ji2023survey}.
Yet standard captioning metrics reward the second output over the first, and a paper reporting only BLEU or BERTScore would conclude that domain adaptation has ``solved'' the problem.

\paragraph{This paper.}
We ask a narrower question than ``can Video-LLMs explain crashes?''.
We ask \emph{which} failures domain adaptation fixes, which it leaves untouched, and whether a popular training-free remedy---decode-time contrastive suppression of temporally implausible tokens~\cite{season2025,vcd}---repairs what remains.
We study this on the Car Crash Dataset (CCD)~\cite{bao2020nondeterministic}, whose 1{,}500 dashcam clips come with structured forensic fields (severity, impact location, vehicles, weather, crash window) and a $\sim$95-word free-text explanation, so that every claim a model makes can be scored against a human-written reference.

We evaluate four open Video-LLMs zero-shot, fine-tune Qwen2.5-VL-7B~\cite{qwen25vl} with QLoRA~\cite{qlora} into \crashlogic{}, and evaluate seven decode-time variants of Temporal Contrastive Decoding (\tcd).
Beyond lexical metrics we score each output for \emph{omission} and \emph{hallucination} with an ARGUS-style dual cost~\cite{argus2025} adapted to CCD's structured fields, for entailment against the forensic ground truth with an NLI cross-encoder, and for temporal grounding against prior-only baselines.
All comparisons use paired tests on aligned per-video scores~\cite{dror2018hitchhiker}.

\paragraph{Findings.}
Three results organise the paper.

\emph{(i) Adaptation removes omissions, not hallucinations.}
Fine-tuning reduces the omission cost from $0.462$ to $0.107$ (improved on 140/150 videos, $p<10^{-4}$) and raises BLEU-4 ninefold.
The hallucination cost is unchanged ($0.227\to0.223$, $p=0.81$).
The adapted model asserts every forensic field for every video, and roughly one field in five is wrong.
On the six test videos whose ground truth states that no collision occurred, the adapted model fabricates one in three cases under greedy decoding and in all six under \tcd.

\emph{(ii) Temporal localisation is a learned prior.}
\crashlogic{} raises mean temporal IoU from $0.012$ to $0.373$---but it emits exactly two crash windows, $[2,3]$\,s or $[3,4]$\,s, and its predicted start time is uncorrelated with the true start (Spearman $\rho{=}0.009$, $p{=}0.91$).
A constant $[3,4]$\,s window scores tIoU $0.408$, higher than any model.
Because frames are supplied without absolute time stamps, the model has no signal from which to localise and instead reproduces the training-set marginal---a shortcut in the sense of Geirhos~\emph{et al.}~\cite{geirhos2020shortcut}.
Reporting tIoU without such a baseline, as prior accident-VLM work does, overstates temporal competence.

\emph{(iii) Temporal contrastive decoding does not repair the residue.}
At $\alpha{=}0.5$, \tcd{} changes no metric significantly (paired Wilcoxon $p\geq0.21$ on all nine metrics; NLI-Score $\Delta{=}{+}0.062$, 95\% CI $[-0.034,+0.156]$).
At $\alpha\geq1.5$ it significantly \emph{increases} hallucination and omission and lowers ROUGE-L.
We give a mechanistic account: a reversed-frame negative preserves every order-invariant cue---vehicle colour, weather, agent identity---so the logit difference cannot discriminate the errors that dominate the residual (wrong agent, wrong severity), and a model that has not learned to use frame order receives almost no contrast signal on timing.

\paragraph{Contributions.}
\begin{itemize}
  \item A held-out, leakage-checked evaluation protocol for dense crash explanation on CCD that scores omission and hallucination separately, includes prior-only temporal baselines, and reports paired statistics over per-video scores.
  \item \crashlogic{}, a QLoRA adaptation of Qwen2.5-VL-7B that is the strongest system on every completeness metric, together with the finding that it leaves hallucination untouched.
  \item A field-level decomposition of hallucination (timestamp, severity, agent colour) with an automatic error taxonomy in the appendix.
  \item A paired evaluation of temporal contrastive decoding on this task, including a negative result with a mechanistic explanation and a demonstration that the collision prior worsens under contrast.
  \item Release of all model outputs, per-video scores, and regeneration scripts.
\end{itemize}

We deliberately place severity classification, structured-field accuracy, and human Likert studies outside the main results: the first two appear only as components of the hallucination analysis, and the third we could not complete at the scale a reviewer should demand (Section~\ref{sec:limitations}).
The system diagram is Figure~\ref{fig:pipeline}.
